\documentclass[12pt,a4paper]{article}
\usepackage[margin=25mm, headheight=15pt, headsep=10pt]{geometry}
\usepackage{fontspec}
\usepackage[table]{xcolor} % Note: table option is required for rowcolors
\usepackage{titlesec}
\usepackage{tocloft}
\usepackage{fancyhdr}
\usepackage{booktabs}
\usepackage{tcolorbox}
\tcbuselibrary{skins, breakable}
\usepackage[colorlinks=true, linkcolor=emerald, urlcolor=emerald, bookmarks=true]{hyperref}

\definecolor{emerald}{RGB}{0,102,51}
\definecolor{darkred}{RGB}{139,0,0}
\definecolor{lightgray}{RGB}{245,245,245}

\usepackage[nil]{babel}
\babelprovide[import, main]{bengali}
\babelprovide[import]{arabic}
\babelfont{rm}[Renderer=HarfBuzz,Scale=1.0]{Noto Serif Bengali}
\babelfont{sf}[Renderer=HarfBuzz]{Noto Sans Bengali}
\babelfont{tt}[Renderer=HarfBuzz]{Noto Sans Bengali}
\babelfont[arabic]{rm}[Renderer=HarfBuzz,Scale=1.1]{Noto Naskh Arabic}

% Section styling
\titleformat{\section}{\Large\bfseries\color{emerald}}{\thesection.}{0.5em}{}
\titleformat{\subsection}{\large\bfseries\color{emerald!85!black}}{\thesubsection.}{0.5em}{}
\titleformat{\subsubsection}{\normalsize\bfseries\color{emerald!70!black}}{\thesubsubsection.}{0.5em}{}
\titlespacing*{\section}{0pt}{18pt}{8pt}

% Fancy Header/Footer setup
\pagestyle{fancy}
\fancyhf{} % clear all header and footer fields
\fancyhead[L]{\color{emerald}\small\textbf{\nouppercase{\leftmark}}}
\fancyfoot[L]{\scriptsize\color{black!50}{\lat{AI}}-সংকলিত, আলেম-যাচাইকৃত নয়}
\fancyfoot[C]{\color{emerald!70!black}\thepage}
\renewcommand{\headrulewidth}{0.4pt}
\renewcommand{\headrule}{\hbox to\headwidth{\color{emerald}\leaders\hrule height \headrulewidth\hfill}}
\renewcommand{\footrulewidth}{0pt}

% Custom boxes
% 1. Ayat/Hadith Box
\newtcolorbox{ayatbox}[1][]{
    enhanced, breakable, 
    colback=emerald!5, 
    colframe=emerald, 
    boxrule=0.5pt, 
    arc=3pt, 
    left=10pt, right=10pt, top=10pt, bottom=10pt,
    #1
}

% 2. Quote/Translation Box
\newtcolorbox{quotebox}[1][]{
    enhanced, breakable,
    frame hidden,
    colback=lightgray,
    borderline west={3pt}{0pt}{emerald},
    left=15pt, right=10pt, top=10pt, bottom=10pt,
    sharp corners,
    #1
}

% 3. AI-generated notice box
\newtcolorbox{warnbox}[1][]{
    enhanced, breakable,
    colback=red!4,
    colframe=darkred,
    boxrule=0.8pt,
    arc=3pt,
    left=10pt, right=10pt, top=10pt, bottom=10pt,
    #1
}

% Commands
\newcommand{\arab}[1]{\foreignlanguage{arabic}{#1}}
\newcommand{\kib}[1]{\textcolor{emerald}{\textbf{#1}}}

% Latin (English) fallback: Noto Serif Bengali has NO Latin glyphs
\newfontfamily\latfont[Renderer=HarfBuzz]{Noto Serif}
\newcommand{\lat}[1]{{\latfont #1}}

\setlength{\parskip}{5pt}
\setlength{\parindent}{0pt}

% Fix for missing bullet in Noto Serif Bengali
\renewcommand{\labelitemi}{$\bullet$}



\begin{document}
\begin{center}{\Large\bfseries\color{emerald} ঘটনা ১০: ইবনে মাসউদ (রাঃ) — পুরনো চাদর পরে অহংকার দমন}\end{center}
\vspace{1em}
\section*{ঘটনা ১০: ইবনে মাসউদ (রাঃ) — পুরনো চাদর পরে অহংকার দমন}

\begin{warnbox}
 \textbf{গুরুত্বপূর্ণ নোটিশ (\lat{Important} \lat{Notice}):} এই কন্টেন্টটি \lat{AI} (কৃত্রিম বুদ্ধিমত্তা) দিয়ে প্রস্তুত ও সংকলিত। \lat{AI} কঠোর সূত্র-মান নীতি মেনে শুধু নির্ভরযোগ্য উৎস থেকে তথ্য সংগ্রহ করেছে — কেবল সহীহ/হাসান হাদিস ও সহীহ সনদের আসার রাখা হয়েছে, দুর্বল সনদ বাদ দেওয়া হয়েছে। তবে এটি কোনো আলেম বা মুহাদ্দিস কর্তৃক যাচাইকৃত নয়।
\end{warnbox}


\medskip\hrule\medskip

\subsection*{১. সূত্র কার্ড}

\begin{table}[h]\centering\small
\begin{tabular}{ll}
বিষয় & বিবরণ \\
\hline
\textbf{বর্ণনাকারী} & আবু মাজলায (রহ.) থেকে \\
\textbf{গ্রন্থ} & মুসান্নাফ ইবনে আবি শায়বাহ, হাদিস ২৬৫২৪ ও ২৫৪০৫ \\
\textbf{অধ্যায়} & মুসান্নাফ — কিতাবুল লিবাস (পোশাক) ও কিতাবুজ-যুহদ \\
\textbf{সনদের মান} & \textbf{সহীহ সনদের আসার} — ইমাম বুখারি হাদিস ৯১-এর রাবিদের কাছ থেকেও বর্ণিত \\
\textbf{মূল বিষয়} & সহীহ হাদিসের প্রভাব সাহাবাদের জীবনে — কিবরের হাদিস শুনে ইবনে মাসউদের আচরণ \\
\hline
\end{tabular}
\end{table}

\medskip\hrule\medskip

\subsection*{২. পটভূমি (কে, কখন, কোথায়, কেন)}

\textbf{আবদুল্লাহ ইবনে মাসউদ (রাঃ)} — ইসলামের প্রসিদ্ধ সাহাবি, নবীজির ঘরের সদস্য, কুরআনের বিশিষ্ট কারী ও আলেম। তিনি কুফায় \lat{teacher} (শিক্ষক) হিসেবে ছিলেন এবং সেখানে তাঁর ছাত্র/সঙ্গীসহ কিছু লোক একসাথে থাকতেন।

শীতের সময় এসেছিল। ইবনে মাসউদের সঙ্গীরা \textbf{ঠান্ডা}য় কষ্ট পেতে শুরু করল। তখন সঙ্গীরা দামি ও গরম \lat{clothes} (পোশাক) পরতে শুরু করল। কিন্তু অদ্ভুত এক সংকোচ তাদের মাঝে ছড়িয়ে পড়ল — কেউ কেউ সস্তা-পুরনো পোশাক পরে আসতে \lat{shame} (লজ্জা) বোধ করল।

ইবনে মাসউদ (রাঃ) এই পরিবেশ বুঝতে পারলেন — মানুষ সস্তা পোশাক পরে আসতে সংকুচিত হচ্ছে, যার অর্থ তাদের মনে কোনো না কোনোভাবে পোশাক নিয়ে গর্ব ও শ্রেষ্ঠত্বের ভাব আছে। তখন তিনি নিজেই একটি পুরনো \lat{cloak} (চাদর) পরে এলেন — এবং তিন দিন পর্যন্ত সেই একই \lat{cloak} পরে এলেন।

\medskip\hrule\medskip

\subsection*{৩. পূর্ণ ঘটনার বর্ণনা}

আবু মাজলায (রহ.) বলেন:

\begin{quote}
\textbf{\lat{“}ইবনে মাসউদের সঙ্গীদেরকে শীত (ঠান্ডা) পীড়িত করল। লোকজন তখন সস্তা ও পুরনো পোশাক (কিসা বা থাওব) পরে আসতে সংকুচিত বোধ করছিল।} \\
\textbf{এরপর আবু আবদির রহমান (ইবনে মাসউদ, রাঃ) একটি চাদরে (আবাহ/আবায়াহ) হাজির হলেন — তারপর দ্বিতীয় দিনও সেই চাদরে, তৃতীয় দিনও সেই একই চাদরে এলেন।} \\
\textbf{এরপর মানুষও ধীরে ধীরে (দামি পোশাক বাদ দিয়ে) সাধারণ পোশাকে ফিরে আসতে শুরু করল।} \\
\textbf{অতঃপর ইবনে মাসউদ (রাঃ) বললেন: "আমি নবীজি (সাল্লাল্লাহু আলাইহি ওয়া সাল্লাম)-কে বলতে শুনেছি: 'যার হৃদয়ে সরিষার দানার পরিমাণ কিবর থাকবে, সে জান্নাতে প্রবেশ করবে না।' তাই আমি আমার কিবর ভাঙতে চেয়েছি।"\lat{”}}
\end{quote}

\medskip\hrule\medskip

\subsection*{৪. শব্দের ব্যাখ্যা}

\begin{table}[h]\centering\small
\begin{tabular}{ll}
শব্দ & অর্থ \\
\hline
\textbf{কিসা (\arab{كساء})} & চাদর/কাপড় — সাধারণ পোশাক \\
\textbf{থাওব (\arab{ثوب})} & জামা/পোশাক \\
\textbf{আবাহ/আবায়াহ (\arab{عباءة})} & মোটা চাদর — সাধারণ, মোটা কাপড়ের চাদর \\
\textbf{দূন (\arab{دون})} & নিম্নমানের/সস্তা \\
\textbf{মিছকাল যাররাহ (\arab{مثقال} \arab{ذرة})} & সরিষার দানার ওজন — অতি সামান্য পরিমাণ \\
\hline
\end{tabular}
\end{table}

\textbf{ইবনে মাসউদের দৃষ্টিভঙ্গি:} তিনি ধনী বা গরিব বলেই পুরনো চাদর পরেননি; তিনি এটি করেছেন \textbf{ইচ্ছাকৃতভাবে} — কিবরের হাদিসের শিক্ষা বাস্তবে প্রয়োগ করতে এবং সঙ্গীদের সংকোচ দূর করতে।

\medskip\hrule\medskip

\subsection*{৫. সংশ্লিষ্ট হাদিস ও আয়াত}

\textbf{হাদিস (মুসলিম ৯১):} নবীজি বলেছেন: \textbf{\lat{“}যার হৃদয়ে সরিষার দানার পরিমাণ কিবর থাকবে, সে জান্নাতে প্রবেশ করবে না।\lat{”}} — এই হাদিসই ইবনে মাসউদের আচরণের ভিত্তি।

\textbf{হাদিস (মুসলিম ৯১বি):} \textbf{\lat{“}যার হৃদয়ে সরিষার দানার পরিমাণ ঈমান থাকবে, সে জাহান্নামে প্রবেশ করবে না; আর যার হৃদয়ে সরিষার দানার পরিমাণ কিবর থাকবে, সে জান্নাতে প্রবেশ করবে না।\lat{”}}

\textbf{হাদিস (বুখারি ৫৭৮৪):} নবীজি বলেছেন: \textbf{\lat{“}যে ব্যক্তি অহংকারে পোশাক টেনে হাঁটে, আল্লাহ কিয়ামতের দিন তার দিকে তাকাবেন না।\lat{”}}

\textbf{আয়াত (সূরা আল-আ'রাফ ৭:৪০):}
\begin{quote}
\textbf{\lat{“}নিশ্চয়ই যারা আমার আয়াতকে মিথ্যা বলেছে ও তা থেকে অহংকার করেছে, তাদের জন্য আকাশের দরজা খোলা হবে না; এবং উট সূঁচের ছিদ্রে প্রবেশ না করা পর্যন্ত তারা জান্নাতে ঢুকবে না।\lat{”}}
\end{quote}

\textbf{আয়াত (সূরা আল-কাসাস ২৮:৮৩):}
\begin{quote}
\textbf{\lat{“}এটি আখেরাতের ঘর, যা আমি দেবো তাদেরকে যারা পৃথিবীতে উচ্চতা (বড়ত্ব) ও বিপর্যয় কামনা করে না।\lat{”}}
\end{quote}

\medskip\hrule\medskip

\subsection*{৬. রেফারেন্স}

\begin{itemize}
\item মুসান্নাফ ইবনে আবি শায়বাহ, হাদিস ২৬৫২৪ ও ২৫৪০৫ (সহীহ সনদ)
\item সহীহ মুসলিম, হাদিস ৯১ (কিবরের হাদিস)
\item সহীহ বুখারি, হাদিস ৫৭৮৪
\item সূরা আল-আ'রাফ, আয়াত ৪০; সূরা আল-কাসাস, আয়াত ৮৩
\end{itemize}

\end{document}
